\section{总结}

在本文中，我们给出了为正在开展的、采用伪 PDF 方法计算胶子 PDF 的工作
奠定基础的结果。

特别地，我们对可能的双胶子关联函数作了分类。
我们指出了其中包含不变振幅 ${\cal M}_{pp}(\nu,-z^2)$ 的那些函数，
该振幅在光锥极限 $z^2 \to 0$ 下决定胶子 PDF。
由于此极限是奇异的，需要联系 ${\cal M}_{pp}(\nu,z_3^2)$ 与光锥 PDF $f(x,\mu^2)$ 的
匹配条件。

为此，我们采用 Ref.~\cite{Balitsky:1987bk} 的方法，
对两个胶子场强张量的规范不变关联器（所有洛伦兹指标均显式保留）
完成了单圈修正的计算。
为保持规范不变性，我们使用了维数正规化。

由于 DR 对与 UV 奇异性相关的对数
和反映 DGLAP 演化的对数给出相同的形式 $\ln z_3^2 \mu^2$，我们尽力分离了
小 $z_3^2$ 处 \mbox{$\ln z_3^2$ 依赖}的这两个来源。
当我们构造约化 ITD ${\mathfrak M}(\nu,z_3^2)$ 时，UV 相关的贡献相消，
约化 ITD 与光锥 ITD 之间的匹配关系中
只剩下 DGLAP 相关的项。

匹配关系也可写成核形式 (\ref{MtchI})，
它直接把 ${\mathfrak M}(\nu,z_3^2)$ 的格点数据
与归一化胶子 PDF $x f_g (x,\mu^2)/ {\langle x \rangle_{\mu^2}}$ 联系起来。
平均胶子动量份额 ${\langle x \rangle_{\mu^2}}$
需要由单独的格点计算提取。

我们还证明，这些结果可用于相当直接地计算
准分布的单圈修正，
为其结构提供新的洞见，
并可用于检验
胶子准分布匹配条件的相关结果。

在后续出版物中，我们计划给出计算的更多细节，
并给出盒图的完整结果，
特别是分布振幅与 GPD 格点计算所需的
非前向运动学。
我们还计划纳入胶子--夸克
与夸克--胶子项的计算。
